\section{Marco teórico}

El marco teórico de cada práctica debe diseñarse bajo los diferentes contenidos abordados 
en el experimento. Por ejemplo, si el contenido temático de la práctica es el diseño de 
la ``Bobina de Tesla'' debe consultar referentes verídicos que expliquen el funcionamiento 
desde la física del experimento~\parencite{cervantes1999}.

Para el caso de las imágenes, tenga en cuenta las normas APA, tal como se presenta
en la Figura 1, además recuerde que cualquier recurso fotográfico deberá citarlo previamente 
en el cuerpo del documento.

Recuerde que debe incluir cada gráfica el nombre de los ejes y el título, además que
debe ser clara y entendible. Se recomienda el uso de software especializado como
Python, Power BI o MATLAB para obtener las gráficas.

\begin{table}[htbp]
\centering
\caption{Representación visual de las geometrías y ecuaciones de la Ley de Biot-Savart.}
\label{tab:biot_visual}

% Definimos columnas: 'l' para el nombre, 'X' para la imagen (se ajusta), 'm{3cm}' para la ecuación centrada verticalmente
\begin{tabularx}{\linewidth}{l >{\centering\arraybackslash}m{5.0cm} X}
\toprule
\textbf{Geometría} & \textbf{Ecuación de $B(r)$} & \textbf{Diagramas} \\
\midrule

% FILA 1: Hilo Recto
Hilo recto infinito & 
% \raisebox ajusta la imagen verticalmente para que el texto de al lado quede centrado
%\raisebox{-.5\height}{\includegraphics[width=0.9\linewidth, height=2cm, keepaspectratio]{images/hilo_recto.png}} & 
$B(r) = \frac{\mu_0 I}{2\pi r}$ &
\raisebox{-.5\height}{\includegraphics[width=0.9\linewidth, height=2cm, keepaspectratio]{imagen1.png}} \\ 
\addlinespace[15pt] % Espacio extra entre filas

% FILA 2: Lazo Circular (Centro)
Espira circular (en el centro) & 
%\raisebox{-.5\height}{\includegraphics[width=0.9\linewidth, height=2cm, keepaspectratio]{images/lazo_circular_centro.png}} & 
$B(r) = \frac{\mu_0 I}{2R}$ &
\raisebox{-.5\height}{\includegraphics[width=0.9\linewidth, height=2cm, keepaspectratio]{imagen1.png}} \\ 
\addlinespace[15pt]

% FILA 3: Solenoide
Solenoide (bobina ideal) & 
%\raisebox{-.5\height}{\includegraphics[width=0.9\linewidth, height=2cm, keepaspectratio]{images/solenoide.png}} & 
$B(r) = \mu_0 \frac{N}{L} I$ &
\raisebox{-.5\height}{\includegraphics[width=0.9\linewidth, height=2cm, keepaspectratio]{imagen1.png}} \\ 
\addlinespace[15pt]

% FILA 4: Lazo Circular (Eje)
Eje del lazo circular & 
%\raisebox{-.5\height}{\includegraphics[width=0.9\linewidth, height=2cm, keepaspectratio]{images/lazo_circular_eje.png}} & 
$B(r) = \frac{\mu_0 I}{2} \frac{R^2}{(x^2 + R^2)^{3/2}}$ &
\raisebox{-.5\height}{\includegraphics[width=0.9\linewidth, height=2cm, keepaspectratio]{imagen1.png}} \\ 
\bottomrule
\end{tabularx}
\end{table}